% !TeX root = elegantnote-en.tex
\section{Condensations and Alpha Substitutions of Carbonyl Componds}

\subsection{Introduction 20260528}

\begin{enumerate}
    \item \textbf{Alpha} ($\boldsymbol{\alpha}$) \textbf{substitutions} replace a hydrogen atom at the $\boldsymbol{\alpha}$ \textbf{carbon atom}  by some other group.
    \item Alpha substitutions generally take place when the carbonyl compound is converted to its enolate ion or its enol tautomer.
\end{enumerate}

\begin{mechanism}
    Alpha Substitution
    \begin{itemize}
        \item \textbf{\textit{Step 1:}} Deprotonation of a carbon to form an enolate.
        \item \textbf{\textit{Step 2:}} Nucleophilic attack on an electrophile.
        \begin{figure}[H]
            \centering
            \includegraphics[width=0.7\textwidth]{22.1_Alpha_Substitution.png}
            %\caption{***}
            \label{22.1_Alpha_Substitution}
        \end{figure}
    \end{itemize}
\end{mechanism}

\begin{enumerate}[resume]
    \item Carbonyl \textbf{condensations}: alpha substitutions where the electrophile is another carbonyl compound.
    \item If the electrophile is a ketone or an aldehyde: nucleophilic addition.
\end{enumerate}

\begin{mechanism}
    Addition of an Enolate to Ketones and Aldehydes (a Condensation)
    \begin{itemize}
        \item \textbf{\textit{Step 1:}} The enolate adds to the carbonyl group.
        \item \textbf{\textit{Step 2:}} Protonation of the alkoxide.
        \begin{figure}[H]
            \centering
            \includegraphics[width=0.7\textwidth]{22.1_Carbonyl_Condensation_1.png}
            %\caption{***}
            \label{22.1_Carbonyl_Condensation_1}
        \end{figure}
    \end{itemize}
\end{mechanism}

\begin{enumerate}[resume]
    \item If the electrophile is an ester: nucleophilic acyl substitution.
\end{enumerate}

\begin{mechanism}
    Substitution of an Enolate on an Ester (a Condensation)
    \begin{itemize}
        \item \textbf{\textit{Step 1:}} Addition of the enolate.
        \item \textbf{\textit{Step 2:}} Elimination of the alkoxide.
        \begin{figure}[H]
            \centering
            \includegraphics[width=0.7\textwidth]{22.1_Carbonyl_Condensation_2.png}
            %\caption{***}
            \label{22.1_Carbonyl_Condensation_2}
        \end{figure}
    \end{itemize}
\end{mechanism}

\begin{enumerate}[resume]
    \item Alpha substitutions and condensations of carbonyl compounds are some of the most common methods for forming carbon-carbon bonds.
    \item Considering the structure and formation of enols and enolate ions.
\end{enumerate}
%----------%